\chapter{Literature and Technology Review}

\section{Introduction}
% TODO: Explain the review scope and how academic, technical, and practitioner sources are used.

% Pattern for drafting from the quote bank:
% 1. Paste the exact quote in a quote environment.
% 2. Add the citation immediately after it.
% 3. Later, rewrite the surrounding discussion in your own words and remove quotes that are no longer needed.

\section{Curvature Preference and Shape Perception}
\subsection{Preference for Curved over Angular Forms}
\begin{quote}
``People show a small but consistent preference for the round over the sharp, for the curvy over the pointy, when judging lines, shapes, abstract objects, and artifacts. This overall preference for curvature is now firmly established.''
\end{quote}
\citep{cotter_curve_2017}
% Later: use this to establish the perceptual/aesthetic basis for investigating curvilinear UI components.

\begin{quote}
``The idea that curved contours appear to be more pleasurable to the human eye than straight lines is a recurrent theme in aesthetic literature that has been occasionally explored from the scientific psychology.''
\end{quote}
\citep{gomez-puerto_preference_2016}
% Later: use this to connect older aesthetic theory with empirical shape-preference research.

\subsection{Affect, Threat, and Processing Fluency}
\begin{quote}
``There was greater bilateral activation in the amygdala for sharp-angled shapes compared with curved shapes.''
\end{quote}
\citep{bar_visual_2007}
% Later: handle this carefully as supporting evidence, not a direct claim that angular UI is threatening.

\subsection{Limitations of Curvature Preference Research}
% TODO: Add quotes that qualify the evidence, e.g. cultural effects, stimulus type, context, and individual differences.

\section{Organic, Biomorphic, and Biophilic Design}
\subsection{Organic Design in Art, Product Design, and Architecture}
% TODO: Add quotes from Agkathidis, Joye, and related design/architecture sources.

\subsection{Biomorphic Forms and Natural Aesthetics}
\begin{quote}
``Biomimetic design uses curves as the aesthetic vehicle, effectively overcoming the mechanical feel of geometric straight lines and imbuing products with organic vitality.''
\end{quote}
\citep{xiao_research_2025}
% Later: use this as a bridge between organic/biomorphic design theory and the project focus on curvilinear interface components.

\subsection{Relevance to Digital Interface Design}
% TODO: Add quotes linking natural/organic aesthetics to digital interaction, engagement, or user experience.

\section{Organic User Interfaces and Shape-Based Interaction}
\subsection{Organic User Interfaces}
\begin{quote}
``With the emergence of flexible display technologies, it will be necessary for interface designers to move beyond flat interfaces and to contextualize interaction in an object's physical shape.''
\end{quote}
\citep{holman_design_2013}
% Later: explain why this project borrows the concern with form and shape, while applying it to web UI rather than flexible displays.

\subsection{Tangible and Non-Planar Interfaces}
% TODO: Add quotes on shape-changing, tangible, or non-planar interaction where relevant.

\subsection{Transferability to Web-Based Interfaces}
% TODO: Critically discuss what can and cannot be transferred from OUI literature to conventional web interfaces.

\section{Procedural Generation and Computational Form}
\subsection{Procedural Generation of Natural Forms}
\begin{quote}
``Organic Information Design is concerned with augmenting one's ability to process large amounts of data [using] simulated organic [processes].''
\end{quote}
\citep{fry_organic_1997}
% Later: use this as historical context for computational organic form, not as a direct web UI implementation source.

\subsection{Randomness, Constraints, and Parameterisation}
% TODO: Add quotes on generative, parametric, or procedural design methods.

\subsection{Generative Tools for Design Assets}
% TODO: Add sources on procedural generation, design tools, and asset generation.

\section{Web UI Layouts, Grids, and Component Systems}
\subsection{Rectilinear Layout Conventions}
% TODO: Add quotes on grids, component systems, Bootstrap/CSS layout, and design systems.

\subsection{CSS, Bootstrap, and Grid-Based Design}
% TODO: Add technical sources that establish the dominant implementation model.

\subsection{Existing Techniques for Curved or Irregular UI Elements}
\begin{quote}
``We present a method for resolution independent rendering of paths and bounded regions, defined by quadratic and cubic spline curves, that leverages the parallelism of programmable graphics hardware to achieve high performance.''
\end{quote}
\citep{loop_resolution_2005}
% Later: use this to show that curve rendering is technically mature, while the dissertation gap is about reusable web UI generation.

\section{Design Science Research and Evaluation Literature}
\subsection{Design Science Research}
\begin{quote}
``Design science, as the other side of the IS research cycle, creates and evaluates IT artifacts intended to solve identified organizational problems.''
\end{quote}
\citep{hevner_design_2004}
% Later: use this to justify treating the API/prototype as the research artefact.

\subsection{Iterative and Exploratory Development}
% TODO: Add quotes on iterative development, exploratory prototyping, versioning, or agile development where useful.

\subsection{Usability and Usefulness Evaluation}
% TODO: Add quotes supporting perceived ease of use, perceived usefulness, SUS, task-based feedback, or developer experience evaluation.

\section{Research and Technology Gap}
% TODO: Synthesise the gap: strong design/perception basis, but limited reusable API support for procedural curvilinear UI components.

\section{Summary}
% TODO: Summarise how the literature informs the prototype requirements and evaluation approach.
